\documentclass[UTF8]{ctexbeamer}

\usepackage{amsmath, amssymb, bm}
\usepackage{hyperref}

\usetheme{Madrid}
\usecolortheme{default}
\usefonttheme{professionalfonts}

\title{Mandelbrot 集简介}
\author{}
\date{}

\begin{document}

\frame{\titlepage}

\begin{frame}{Mandelbrot 集是什么？}
\begin{itemize}
    \item Mandelbrot 集（曼德勃罗集）是复动力系统中的一个经典分形对象。
    \item 它由如下迭代过程定义：
    \[
        z_{n+1} = z_n^2 + c
    \]
    其中：
    \begin{itemize}
        \item \(c \in \mathbb{C}\) 是一个固定的复数参数；
        \item \(z_0 = 0\)；
        \item 然后不断迭代生成 \(z_1, z_2, z_3, \dots\)
    \end{itemize}
    \item 对每一个复数 \(c\)，观察这个数列是否始终有界。
\end{itemize}

\medskip
\textbf{定义：}
\[
M = \left\{ c \in \mathbb{C} \;\middle|\; \{z_n\}_{n=0}^{\infty} \text{ 在 } z_0=0,\ z_{n+1}=z_n^2+c \text{ 下有界} \right\}
\]
这就是 \textbf{Mandelbrot 集}。
\end{frame}

\begin{frame}{为什么它是“集合”？}
\begin{itemize}
    \item 这里真正变化的不是 \(z_0\)，而是参数 \(c\)。
    \item 每取一个复数 \(c\)，就得到一条轨道：
    \[
    0,\quad c,\quad c^2+c,\quad (c^2+c)^2+c,\quad \dots
    \]
    \item 如果这条轨道不会无限发散，那么这个 \(c\) 就属于 Mandelbrot 集。
    \item 因此，Mandelbrot 集本质上是复平面上的一个点集：
    \begin{itemize}
        \item 横轴表示 \(\operatorname{Re}(c)\)；
        \item 纵轴表示 \(\operatorname{Im}(c)\)。
    \end{itemize}
\end{itemize}
\end{frame}

\begin{frame}{生成 Mandelbrot 集的核心思想}
对复平面上的每一个点 \(c=x+iy\)，执行以下步骤：

\begin{enumerate}
    \item 令初值 \(z_0 = 0\)；
    \item 按照递推公式
    \[
        z_{n+1}=z_n^2+c
    \]
    反复计算；
    \item 在计算过程中检查 \(|z_n|\) 是否超过阈值；
    \item 若很快超过阈值，则认为该点 \(c\) 不属于 Mandelbrot 集；
    \item 若迭代很多次后仍未超过阈值，则认为该点大概率属于 Mandelbrot 集。
\end{enumerate}
\end{frame}

\begin{frame}{为什么只需检查 \texorpdfstring{\(|z_n|>2\)}{|z_n| > 2}？}
对于迭代
\[
z_{n+1}=z_n^2+c,
\]
有一个重要结论：

\medskip
\textbf{逃逸判据：}
如果某一步出现
\[
|z_n|>2,
\]
那么后续迭代将趋于无穷大。

\medskip
因此：
\begin{itemize}
    \item 一旦 \(|z_n|>2\)，就可以立刻停止；
    \item 该参数 \(c\) 不属于 Mandelbrot 集。
\end{itemize}
\end{frame}

\begin{frame}{详细生成规则（数学版）}
设要绘制的复平面区域为
\[
x \in [x_{\min}, x_{\max}], \qquad
y \in [y_{\min}, y_{\max}].
\]

对每个像素点 \((i,j)\)，做如下映射：
\[
c = x + iy,
\]
其中
\[
x = x_{\min} + \frac{i}{W-1}(x_{\max}-x_{\min}),
\qquad
y = y_{\min} + \frac{j}{H-1}(y_{\max}-y_{\min}).
\]

然后定义：
\[
z_0 = 0,
\qquad
z_{n+1}=z_n^2+c.
\]

给定最大迭代次数 \(N_{\max}\)：
\begin{itemize}
    \item 若存在最小整数 \(n\le N_{\max}\)，使得 \(|z_n|>2\)，则该点逃逸；
    \item 若对所有 \(0\le n\le N_{\max}\) 都有 \(|z_n|\le 2\)，则把该点视为集合内部点。
\end{itemize}
\end{frame}

\begin{frame}{详细生成规则（算法版）}
\textbf{输入：}
\begin{itemize}
    \item 图像宽度 \(W\)、高度 \(H\)
    \item 复平面窗口 \([x_{\min},x_{\max}] \times [y_{\min},y_{\max}]\)
    \item 最大迭代次数 \(N_{\max}\)
\end{itemize}

\textbf{对每个像素：}
\begin{enumerate}
    \item 由像素坐标计算复数 \(c\)；
    \item 初始化 \(z \leftarrow 0\)，\(n \leftarrow 0\)；
    \item 当 \(n < N_{\max}\) 且 \(|z| \le 2\) 时：
    \[
    z \leftarrow z^2 + c,\qquad n \leftarrow n+1
    \]
    \item 记录该点的迭代次数 \(n\)。
\end{enumerate}

\textbf{输出：}
\begin{itemize}
    \item 若 \(n=N_{\max}\)，通常画成黑色；
    \item 若 \(n<N_{\max}\)，按 \(n\) 的大小着色。
\end{itemize}
\end{frame}

\begin{frame}{颜色是如何确定的？}
设某点在第 \(n\) 次迭代第一次满足
\[
|z_n|>2.
\]
则可以用 \(n\) 决定颜色：

\begin{itemize}
    \item \(n\) 小：逃逸快；
    \item \(n\) 大：更接近边界；
    \item \(n=N_{\max}\)：通常视为未逃逸，画成黑色。
\end{itemize}

\medskip
因此，Mandelbrot 集内部通常显示为黑色，外部根据逃逸次数形成彩色层次。
\end{frame}

\begin{frame}{一个具体例子}
考虑参数
\[
c=1.
\]
从 \(z_0=0\) 开始迭代：
\[
z_1=1,\qquad z_2=2,\qquad z_3=5.
\]
于是
\[
|z_3|=5>2,
\]
所以轨道发散，故
\[
1 \notin M.
\]

\medskip
再考虑
\[
c=0.
\]
则
\[
z_0=0,\quad z_1=0,\quad z_2=0,\quad \dots
\]
始终有界，所以
\[
0 \in M.
\]
\end{frame}

\begin{frame}{为什么边界如此复杂？}
\begin{itemize}
    \item Mandelbrot 集由非常简单的迭代规则定义：
    \[
    z_{n+1}=z_n^2+c
    \]
    \item 但在参数空间中，它会产生极其复杂的边界；
    \item 边界附近对参数 \(c\) 的微小变化非常敏感；
    \item 这正是分形几何中“简单规则产生复杂结构”的典型现象。
\end{itemize}
\end{frame}

\begin{frame}[fragile]{伪代码}
\small
\begin{block}{Mandelbrot 集生成伪代码}
\begin{verbatim}
for each pixel (i, j):
    c = map_pixel_to_complex(i, j)
    z = 0
    n = 0

    while n < Nmax and |z| <= 2:
        z = z*z + c
        n = n + 1

    if n == Nmax:
        color = black
    else:
        color = palette(n)
\end{verbatim}
\end{block}

\begin{itemize}
    \item \texttt{map\_pixel\_to\_complex}：把像素映射到复平面；
    \item \texttt{palette(n)}：根据逃逸次数选择颜色。
\end{itemize}
\end{frame}

\begin{frame}{总结}
\begin{itemize}
    \item Mandelbrot 集由
    \[
    z_{n+1}=z_n^2+c,\qquad z_0=0
    \]
    定义；
    \item 参数 \(c\) 是否属于集合，取决于对应轨道是否有界；
    \item 数值绘制时，使用逃逸判据 \(|z_n|>2\)；
    \item 实际绘图流程包括：像素映射、迭代计算、逃逸判断和按次数着色。
\end{itemize}
\end{frame}

\end{document}